\section*{Capítulo \ref{ch:g12:buffers-predominance} --- Disoluciones tampón y diagramas de predominio}

\begin{solution}{exo:g12:buffers-predominance:1}
\omterm{def:g9:acids-bases-ph:ph}{pH} 2: \ce{CH3COOH}. \omterm{def:g9:acids-bases-ph:ph}{pH} 4.76: ambas en cantidades iguales. \omterm{def:g9:acids-bases-ph:ph}{pH} 7:
\ce{CH3COO-}.
\end{solution}

\begin{solution}{exo:g12:buffers-predominance:2}
\omterm{def:g9:acids-bases-ph:ph}{pH} 4: $1/(1 + 10^{0.76}) = 0.15$. \omterm{def:g9:acids-bases-ph:ph}{pH} 6: $1/(1 + 10^{-1.24}) = 0.95$.
\end{solution}

\begin{solution}{exo:g12:buffers-predominance:3}
Amarillo; verde (dentro de su \omterm{def:g12:buffers-predominance:indicator}{zona de viraje}); azul.
\end{solution}

\begin{solution}{exo:g12:buffers-predominance:4}
A \omterm{def:g9:acids-bases-ph:ph}{pH} 7.4, por debajo de 9.26: \ce{NH4+}. A \omterm{def:g9:acids-bases-ph:ph}{pH} 11: \ce{NH3}.
\end{solution}

\begin{solution}{exo:g12:buffers-predominance:5}
Un indicador es a su vez un \omterm{def:g12:ka-and-pka:bronsted}{par ácido--base}: en grandes cantidades
cambiaría el \omterm{def:g9:acids-bases-ph:ph}{pH} que se supone que debe mostrar.
\end{solution}

\begin{solution}{exo:g12:buffers-predominance:6}
$10^{\mathrm{pH} - 4.76}$: 0.1; 1; 10.
\end{solution}

\begin{solution}{exo:g12:buffers-predominance:7}
La fenolftaleína (8.0--10.0) o el azul de timol (8.0--9.1): sus \omterm{def:g12:buffers-predominance:indicator}{zonas de
viraje} contienen 8.7. El rojo de metilo vira entre 4.2 y 6.3, lejos de
8.7: habría cambiado mucho antes.
\end{solution}

\begin{solution}{exo:g12:buffers-predominance:8}
$\mathrm{pH} = 4.76 + \log(0.10/0.20) = 4.76 - 0.30 = 4.46$.
\end{solution}

\begin{solution}{exo:g12:buffers-predominance:9}
\omterm{def:g9:acids-bases-ph:ph}{pH} 1: el \omterm{def:g9:ions:ion}{catión} \ce{H3N+-CH2-COOH}, carga $+1$. \omterm{def:g9:acids-bases-ph:ph}{pH} 6: el zwitterión,
carga total 0. \omterm{def:g9:acids-bases-ph:ph}{pH} 12: el \omterm{def:g9:ions:ion}{anión} \ce{H2N-CH2-COO-}, carga $-1$.
\end{solution}

\begin{solution}{exo:g12:buffers-predominance:10}
Agua: de 7.0 a $-\log(\num{1.0e-4}/0.101) = 3.0$, una bajada de 4
unidades. Tampón: de 4.76 a $4.76 + \log(9.9/10.1) = 4.75$, una bajada de
0.01.
\end{solution}

\begin{solution}{exo:g12:buffers-predominance:11}
\ce{NH4+}/\ce{NH3} ($pK_a = 9.26$). Razón
$[\ce{NH3}]/[\ce{NH4+}] = 10^{9.5 - 9.26} = 1.7$.
\end{solution}

\begin{solution}{exo:g12:buffers-predominance:12}
El \omterm{def:g9:acids-bases-ph:ph}{pH} baja una unidad cuando la razón llega a $1/10$:
$(10 - n)/(10 + n) = 0.1$, así que $n = \qty{8.2}{mmol}$. Un tampón que
contiene más de cada forma puede absorber más ácido o más base para el
mismo cambio de la razón: tiene una capacidad mayor.
\end{solution}

\begin{solution}{exo:g12:buffers-predominance:13}
Razón $10^{5.0 - 4.76} = 1.74$; el ácido suma \qty{10}{mmol}, así que
hacen falta \qty{17.4}{mmol} de etanoato, es decir, \qty{174}{mL} de la
disolución.
\end{solution}

\begin{solution}{exo:g12:buffers-predominance:14}
Tiene dos \omterm{def:g12:ka-and-pka:bronsted}{pares ácido--base}, con dos $pK_a$, y por tanto dos cambios de
color: rojo a \omterm{def:g9:acids-bases-ph:ph}{pH} 1, amarillo a \omterm{def:g9:acids-bases-ph:ph}{pH} 5, azul a \omterm{def:g9:acids-bases-ph:ph}{pH} 10.
\end{solution}

\begin{solution}{exo:g12:buffers-predominance:15}
La \omterm{def:g10:concentration-and-dilution:dilution}{dilución} divide las dos concentraciones por 10: la razón, y por tanto
el \omterm{def:g9:acids-bases-ph:ph}{pH} (4.76), no cambian. Un \omterm{def:g12:ka-and-pka:strong-acid}{ácido fuerte} diluido diez veces sube una
unidad de \omterm{def:g9:acids-bases-ph:ph}{pH}.
\end{solution}

\begin{solution}{pb:g12:buffers-predominance:1}
\textbf{1.} \ce{CO2 + 2H2O <=> HCO3- + H3O+}.

\textbf{2.} El dióxido de carbono disuelto (con el agua) es el ácido; el
hidrogenocarbonato, la base.

\textbf{3.} El valor de la tabla corresponde a \qty{25}{\celsius} en agua
pura; la sangre está a \qty{37}{\celsius} y llena de otros \omterm{def:g9:ions:ion}{iones}, que
modifican la constante (y los fisiólogos cuentan todo el dióxido de
carbono disuelto).

\textbf{4.} A \qty{24}{mmol/L}, $0.024 \times 5.0 = \qty{0.12}{mol}$.

\textbf{5.} Un eje de \omterm{def:g9:acids-bases-ph:ph}{pH} cortado en 6.1: \ce{CO2} por debajo, \ce{HCO3-}
por encima.

\textbf{6.} A 7.4, por encima de 6.1: el hidrogenocarbonato.

\textbf{7.} Ligeramente básica: su \omterm{def:g9:acids-bases-ph:ph}{pH} es superior a 7.

\textbf{8.} $10^{7.4 - 6.1} = 10^{1.3} = 20$.

\textbf{9.} $24 / 20 = \qty{1.2}{mmol/L}$.

\textbf{10.} $10^{1.28} = 19$ y $10^{1.32} = 21$.

\textbf{11.} $20/21 = \qty{95}{\percent}$.

\textbf{12.} $10^{7.6 - 6.1} = 10^{1.5} = 32$.

\textbf{13.} El \omterm{def:g9:ions:ion}{ion} hidrogenocarbonato:
\ce{HCO3- + H3O+ -> CO2 + 2H2O}.

\textbf{14.} $10^{7.1 - 6.1} = 10$.

\textbf{15.} De $24/20 = 1.2$ a $24/10 = \qty{2.4}{mmol/L}$: se duplica.
Los pulmones respiran más deprisa y más hondo, para expulsar el exceso de
dióxido de carbono.

\textbf{16.} Un \omterm{def:g12:ka-and-pka:strong-acid}{ácido débil} y su base pueden absorber, cada uno, lo que
se añade, y el \omterm{def:g9:acids-bases-ph:ph}{pH} solo depende de su razón; un \omterm{def:g12:ka-and-pka:strong-acid}{ácido fuerte} solo
empujaría el \omterm{def:g9:acids-bases-ph:ph}{pH} en un sentido.

\textbf{17.} Nada: la razón, y por tanto el \omterm{def:g9:acids-bases-ph:ph}{pH}, siguen siendo los mismos.

\textbf{18.} Alrededor de 20.
\end{solution}
